\section{Introducción: abrir la caja negra del Coding Agent}

En esta entrega se continúa explorando Codex; después de las anteriores centradas en Context Engineering, Runtime y Harness, esta se enfoca en el diseño de las \textbf{12 herramientas predefinidas} (Tools) internas de Codex. La idea central que el docente enfatiza en repetidas ocasiones es: \textbf{no hay magia, nunca trates al Agent como una caja negra}: hay que revisar y guiar el procesamiento interno del Agent, sobre todo cuando el modelo cae en un ciclo infinito o su comportamiento no corresponde a lo esperado.

\begin{knowledgebox}{Las 12 Tools predefinidas de Codex}
\begin{enumerate}
    \item Apply Patch (editar archivos de código)
    \item Execute Commands (ejecutar comandos)
    \item Write Stdin (escritura en la entrada estándar)
    \item Update Plan (actualizar el plan)
    \item Request User Input (solicitar entrada del usuario)
    \item Web Search (búsqueda web)
    \item View Image (ver imagen)
    \item Spawn Agents (crear sub-Agents)
    \item Send Inputs (enviar instrucciones al sub-Agent)
    \item Wait (esperar al sub-Agent)
    \item Close Agents (cerrar sub-Agents)
    \item Resume (reanudar sub-Agent)
\end{enumerate}
\end{knowledgebox}

